\subsection{SL(2, k) 的一些元素}

对任何 \(t\in k^{*}\) ，置

\[
\begin{array}{l} \theta (t) = e ^ {t X _ {+}} e ^ {t ^ {- 1} X _ {-}} e ^ {t X _ {+}} \\ \qquad = \left( \begin{array}{c c} 1 & t \\ 0 & 1 \end{array} \right) \left( \begin{array}{c c} 1 & 0 \\ - t ^ {- 1} & 1 \end{array} \right) \left( \begin{array}{c c} 1 & t \\ 0 & 1 \end{array} \right) \\ \qquad = \left( \begin{array}{c c} 0 & t \\ - t ^ {- 1} & 0 \end{array} \right) \\ \qquad = e ^ {t ^ {- 1} X _ {-}} e ^ {t X _ {+}} e ^ {t ^ {- 1} X _ {-}}. \end{array}
\]

用第 3 节的记号，

\[
\theta (t) u = - t ^ {- 1} v \quad \theta (t) v = t u
\]

于是

\[
\theta (t) e _ {n} ^ {(m)} = (- 1) ^ {m - n} t ^ {2 n - m} e _ {m - n} ^ {(m)}.\tag{4}
\]

于是元素 \(\theta(t)^2 = \left( \begin{array}{cc} -1 & 0 \\ 0 & -1 \end{array} \right)\) 以 \((-1)^m\) 在 \(\mathrm{V}(m)\) 上作用。若 \(\mathrm{E}\) 是奇维单 \(\mathfrak{sl}(2,k)\) -模，则 \(\theta(t)_{\mathrm{E}}\) 是向量空间 \(\mathrm{E}\) 的对合自同构。特别地，取 \(\mathrm{E}\) 为伴随表示：

\[
\theta (t) _ {\mathrm{E}} X _ {+} = t ^ {- 2} X _ {-} \quad \theta (t) _ {\mathrm{E}} X _ {-} = t ^ {2} X _ {+} \quad \theta (t) _ {\mathrm{E}} H = - H\tag{5}
\]

于是 \(\theta(1)_{\mathrm{E}} = \theta(-1)_{\mathrm{E}}\) 是 \(\mathfrak{sl}(2,k)\) 的典则对合。

对任何 \(t\in k^{*}\) ，置

\[
h (t) = \left( \begin{array}{c c} t & 0 \\ 0 & t ^ {- 1} \end{array} \right) = \theta (t) \theta (- 1).
\]

则 \(h(t)u = tu\) ，\(h(t)v = t^{-1}v\) ，于是

\[
h (t) e _ {n} ^ {(m)} = t ^ {m - 2 n} e _ {n} ^ {(m)}.\tag{6}
\]

命题 6. 设 E 是有限维 \(\mathfrak{sl}(2,k)\) -模，且 \(t\in k^{*}\) 。设 \(\mathrm{E}_p\) 是 E 中权为 \(p\) 的元素之集。

\begin{enumerate}
\def\labelenumi{(\roman{enumi})}
\item
  \(\theta(t)_{\mathrm{E}}|\mathrm{E}_{p}\) 是从 \(\mathrm{E}_p\) 到 \(\mathrm{E}_{-p}\) 的双射。
\item
  \(h(t)_{\mathrm{E}}|\mathrm{E}_p\) 是比为 \(t^p\) 的位似（在 \(\mathrm{E}_p\) 上的）。
\end{enumerate}

若 \(\mathrm{E} = \mathrm{V}(n)\) ，本命题由公式 (4) 与 (6) 得出。一般情形由定理 1 得出。

推论. 设 \(E = E' \oplus E''\) 是命题 2 的推论中定义的 E 的分解。元素 \(\begin{pmatrix}-1 & 0 \\ 0 & -1\end{pmatrix}\) （\(\mathbf{SL}(2, k)\) 的）在 \(E'\) 上以 +1 作用，在 \(E''\) 上以 -1 作用。

这由应用于 \(t = -1\) 的 (ii) 得出。

